\section{Sufficiency}\label{sec:sufficiency}
Sufficiency is a conjunction of floors on separately reported factors; it is recorded, never scored as a probability.
\begin{itemize}
\item \textbf{S0}: a label for ``an S1 floor fails''; it appears only inside \INS.
\item \textbf{S1}: at least one deterministic-confirmed \texttt{same} input valid on both versions (not lossy and not degenerate); all changed units hit; no quarantined input that may hit a changed unit; no truncation, blind spot or tamper flag.
\item \textbf{S2}: S1 plus branch-direction coverage of changed regions of at least the floor (default 0.80) per version, every contributing input deterministic-confirmed, and no lossy-equal observation on return or side channels.
\item \textbf{S3}: specified (plateau, trace stability, second explorer) but \textbf{not implemented}: it is refused in configuration and never emitted. The S3 row of the risk--coverage table (Table~\ref{tab:riskcov}) therefore holds only witnessed divergences.
\end{itemize}
The 0.80 floor is an unvalidated default and no calibration was performed. Experiment E5 varied it (Section~\ref{sec:ablation}). Branch coverage counts conditional jumps, loop and send instructions only, with a conservative denominator. An independent review showed that S2 can be reached from two inputs, because levels do not depend on input count or diversity. S1 and S2 must therefore not be read as a quality grade or as ``80\% of changed branch directions validated as behaviourally safe'', and no evidence shows that they discriminate correct from incorrect verdicts. Sufficiency is also a contested proxy~\cite{inozemtseva2014coverage,just2014mut}.
